%%% Переопределение именований, если иначе не сработает %%%
\gappto\captionsrussian{
%    \renewcommand{\chaptername}{Глава}
   \renewcommand{\appendixname}{Приложение} % (ГОСТ Р 7.0.11-2011, 5.7)
}

%%% Изображения %%%
\graphicspath{{images/}{Dissertation/images/}}         % Пути к изображениям

% Use \MakeTextUppercase (preserves hyphenation) instead of \MakeUppercase
% which disables hyphenation and can break justification in the ToC.
\RequirePackage{textcase}

%%% Интервалы %%%
%% По ГОСТ Р 7.0.11-2011, пункту 5.3.6 требуется полуторный интервал
%% Реализация средствами класса (на основе setspace) ближе к типографской классике.
%% И правит сразу и в таблицах (если со звёздочкой)
% \DoubleSpacing*     % Двойной интервал
\OnehalfSpacing*    % Полуторный интервал
% \SingleSpacing*     % Одиночный интервал
% Отделять формулы от текста сверху и снизу одним строковым интервалом
% (устанавливаем величины вокруг display-math в начале документа)
\AtBeginDocument{%
    \setlength{\abovedisplayskip}{\baselineskip}%
    \setlength{\belowdisplayskip}{\baselineskip}%
    \setlength{\abovedisplayshortskip}{\baselineskip}%
    \setlength{\belowdisplayshortskip}{\baselineskip}%
}
%\setSpacing{1.42}   % Полуторный интервал, подобный Ворду (возможно, стоит включать вместе с предыдущей строкой)

%%% Макет страницы %%%
% Выставляем значения полей (ГОСТ 7.0.11-2011, 5.3.7)
\geometry{a4paper, top=2cm, bottom=2cm, left=2.5cm, right=1cm, nofoot, nomarginpar} %, heightrounded, showframe
\setlength{\topskip}{0pt}   %размер дополнительного верхнего поля
\setlength{\footskip}{12.3pt} % снимет warning, согласно https://tex.stackexchange.com/a/334346

%%% Выравнивание и переносы %%%
%% http://tex.stackexchange.com/questions/241343/what-is-the-meaning-of-fussy-sloppy-emergencystretch-tolerance-hbadness
%% http://www.latex-community.org/forum/viewtopic.php?p=70342#p70342
\tolerance 1414
\hbadness 1414
\emergencystretch 1.5em % В случае проблем регулировать в первую очередь
\hfuzz 0.3pt
\vfuzz \hfuzz
%\raggedbottom
\sloppy                 % Избавляемся от переполнений
% Отключение переносов по всему документу
\hyphenpenalty=10000
\exhyphenpenalty=10000
\clubpenalty=10000      % Запрещаем разрыв страницы после первой строки абзаца
\widowpenalty=10000     % Запрещаем разрыв страницы после последней строки абзаца
\brokenpenalty=4991     % Ограничение на разрыв страницы, если строка заканчивается переносом

\renewcommand{\theintvl}{2} % Интервал между заголовками и текстом в единицах размера шрифта (ГОСТ Р 7.0.11-2011, 5.3.5)
%%% Блок управления параметрами для выравнивания заголовков в тексте %%%
\setlength{\parindent}{1.25em}
\newlength{\otstuplen}
\setlength{\otstuplen}{\theotstup\parindent}
\ifnumequal{\value{headingalign}}{0}{% выравнивание заголовков в тексте
    \newcommand{\hdngalign}{\centering}                % по центру
    \newcommand{\hdngaligni}{}% по центру
    \setlength{\otstuplen}{0pt}
}{%
    \newcommand{\hdngalign}{}                 % по левому краю
    \newcommand{\hdngaligni}{\hspace{\otstuplen}}      % по левому краю
} % В обоих случаях вроде бы без переноса, как и надо (ГОСТ Р 7.0.11-2011, 5.3.5)

\setlength{\cftchapterindent}{0pt}
%%% Оглавление %%%
\setpnumwidth{1.6em}
\renewcommand{\cftdotsep}{0.1}
\renewcommand{\cftchapterdotsep}{\cftdotsep}                % отбивка точками до номера страницы начала главы/раздела

%% Переносить слова в заголовке не допускается (ГОСТ Р 7.0.11-2011, 5.3.5). Заголовки в оглавлении должны точно повторять заголовки в тексте (ГОСТ Р 7.0.11-2011, 5.2.3). Прямого указания на запрет переносов в оглавлении нет, но по той же логике невнесения искажений в смысл, лучше в оглавлении не переносить:
% By default the template used a right-margin for ToC titles which forces
% ragged-right typesetting (no hyphenation). For pdflatex we prefer full
% justification in the ToC so hyphenation can work — set the right margin to 0.
\setrmarg{0pt}
\renewcommand{\cftchapterpagefont}{\normalfont}        % нежирные номера страниц у глав в оглавлении
\renewcommand{\cftchapterleader}{\cftdotfill{\cftchapterdotsep}}% нежирные точки до номеров страниц у глав в оглавлении
% Сделать названия глав в оглавлении небold и в верхнем регистре
\renewcommand{\cftchapterfont}{\normalfont\MakeTextUppercase} % нежирные названия глав в оглавлении


\ifnumgreater{\value{headingdelim}}{0}{%
    % \renewcommand\cftchapteraftersnum{.\space}       % добавляет точку с пробелом после номера раздела в оглавлении
}{}

% Всегда ставим точку после номера для секций (в тексте и в оглавлении)
% \renewcommand\cftsectionaftersnum{.\space}       % добавляет точку с пробелом после номера подраздела в оглавлении
% \renewcommand\cftsubsectionaftersnum{.\space}    % добавляет точку с пробелом после номера подподраздела в оглавлении
% \renewcommand\cftsubsubsectionaftersnum{.\space} % добавляет точку с пробелом после номера подподподраздела в оглавлении

% В текстовых заголовках: всегда ставим точку и пробел после номера
\AfterEndPreamble{%
    \setsecnumformat{\csname the#1\endcsname\space}%
}

% \renewcommand*{\cftappendixname}{\appendixname\space} % Слово Приложение в оглавлении

%%% Колонтитулы %%%
% Порядковый номер страницы печатают на середине верхнего поля страницы (ГОСТ Р 7.0.11-2011, 5.3.8)
\makeevenhead{plain}{}{\rmfamily\thepage}{}
\makeoddhead{plain}{}{\rmfamily\thepage}{}
\makeevenfoot{plain}{}{}{}
\makeoddfoot{plain}{}{}{}
\pagestyle{plain}

% Один интервал сверху и снизу от рисунков и таблиц (ГОСТ) %%%
% \setlength{\intextsep}{\baselineskip}      % над и под float'ами в середине страницы
% \setlength{\textfloatsep}{\baselineskip}   % между текстом и float'ом вверху/внизу страницы

%%% добавить Стр. над номерами страниц в оглавлении
%%% http://tex.stackexchange.com/a/306950
\newif\ifendTOC

\newcommand*{\tocheader}{
%\ifnumequal{\value{pgnum}}{1}{%
%    \ifendTOC\else\hbox to \linewidth%
%      {\noindent{}~\hfill{Стр.}}\par%
%      \ifnumless{\value{page}}{3}{}{%
%        \vspace{0.5\onelineskip}
%      }
%      \afterpage{\tocheader}
%    \fi%
%}{}%
}%

%%% Оформление заголовков глав, разделов, подразделов %%%
%% Работа должна быть выполнена ... размером шрифта 12-14 пунктов (ГОСТ Р 7.0.11-2011, 5.3.8). То есть не должно быть надписей шрифтом более 14. Так и поставим.
%% Эти установки будут давать одинаковый результат независимо от выбора базовым шрифтом 12 пт или 14 пт
\newcommand{\basegostsectionfont}{\fontsize{14pt}{14pt}\selectfont\bfseries}


\makeatletter
\makechapterstyle{thesisgost}{%
    \chapterstyle{default}
    \setlength{\beforechapskip}{0pt}
    \setlength{\midchapskip}{0pt}
    \setlength{\afterchapskip}{\theintvl\curtextsize}
    \renewcommand*{\chapnamefont}{\basegostsectionfont}
    \renewcommand*{\chapnumfont}{\basegostsectionfont}
    \renewcommand*{\chaptitlefont}{\basegostsectionfont}
    \renewcommand*{\chapterheadstart}{}
    % \renewcommand*{\afterchapternum}{.\space}   % добавляет точку с пробелом после номера раздела
    % \ifnumgreater{\value{headingdelim}}{0}{%
    %     \renewcommand*{\afterchapternum}{.\space}   % добавляет точку с пробелом после номера раздела
    % }{%
    %     \renewcommand*{\afterchapternum}{\quad}     % добавляет \quad после номера раздела
    %     % \renewcommand*{\afterchapternum}{.\space}   % добавляет точку с пробелом после номера раздела
    % }
    % Не выводим имя и номер отдельно — соберём "Глава N." вместе с заголовком
    \renewcommand*{\printchapternum}{}%
    \renewcommand*{\printchaptername}{}%
    \renewcommand*{\printchapternonum}{\hdngaligni\hdngalign}%
    % Печать: "Глава N. ЗАГОЛОВОК" в верхнем регистре только для нумерованных глав.
    % Для ненумерованных (\chapter*) — только заголовок в верхнем регистре.
    \renewcommand*{\printchaptertitle}[1]{%
        \hdngaligni\hdngalign
        \chaptitlefont
        \ifthischapternumbered
            \MakeTextUppercase{\@chapapp\space\thechapter\space##1}%
        \else
            \MakeTextUppercase{##1}%
        \fi
    }%
}
\makeatother

% Заменяем предыдущую (нерабочую) обёртку \chapter на надёжную установку флага
\makeatletter
\newif\ifthischapternumbered
% Сохраним оригинальные реализации
\let\orig@makechapterhead\@makechapterhead
\let\orig@makeschapterhead\@makeschapterhead
% Для нумерованных глав: пометить и вызвать оригинал
\renewcommand{\@makechapterhead}[1]{%
  \thischapternumberedtrue
  \orig@makechapterhead{#1}%
}
% Для ненумерованных (\chapter*): пометить как ненумерованные и вызвать оригинал
\renewcommand{\@makeschapterhead}[1]{%
  \thischapternumberedfalse
  \orig@makeschapterhead{#1}%
}
\makeatother

\chapterstyle{thesisgost}

\setsecheadstyle{\basegostsectionfont\hdngalign}
\setsecindent{\otstuplen}

\setsubsecheadstyle{\basegostsectionfont\hdngalign}
\setsubsecindent{\otstuplen}

\setsubsubsecheadstyle{\basegostsectionfont\hdngalign}
\setsubsubsecindent{\otstuplen}

\sethangfrom{\noindent #1} %все заголовки подразделов центрируются с учетом номера, как block

\ifnumequal{\value{chapstyle}}{1}{%
    \chapterstyle{thesisgostchapname}
    \renewcommand*{\cftchaptername}{\chaptername\space} % будет вписано слово Глава перед каждым номером раздела в оглавлении
}{}%

%%% Интервалы между заголовками
\setbeforesecskip{\theintvl\curtextsize}% Заголовки отделяют от текста сверху и снизу тремя интервалами (ГОСТ Р 7.0.11-2011, 5.3.5).
\setaftersecskip{\theintvl\curtextsize}
\setbeforesubsecskip{\theintvl\curtextsize}
\setaftersubsecskip{\theintvl\curtextsize}
\setbeforesubsubsecskip{\theintvl\curtextsize}
\setaftersubsubsecskip{\theintvl\curtextsize}

%%% Вертикальные интервалы глав (\chapter) в оглавлении как и у заголовков
\setlength{\cftbeforechapterskip}{0pt plus 0pt}   % Убираем лишние интервалы между главами в оглавлении


%%% Блок дополнительного управления размерами заголовков
\ifnumequal{\value{headingsize}}{1}{% Пропорциональные заголовки и базовый шрифт 14 пт
    \renewcommand{\basegostsectionfont}{\large\bfseries}
    \renewcommand*{\chapnamefont}{\Large\bfseries}
    \renewcommand*{\chapnumfont}{\Large\bfseries}
    \renewcommand*{\chaptitlefont}{\Large\bfseries}
}{}

%%% Счётчики %%%

%% Упрощённые настройки шаблона диссертации: нумерация формул, таблиц, рисунков
\ifnumequal{\value{contnumeq}}{1}{%
    \counterwithout{equation}{chapter} % Убираем связанность номера формулы с номером главы/раздела
}{}
\ifnumequal{\value{contnumfig}}{1}{%
    \counterwithout{figure}{chapter}   % Убираем связанность номера рисунка с номером главы/раздела
}{}
\ifnumequal{\value{contnumtab}}{1}{%
    \counterwithout{table}{chapter}    % Убираем связанность номера таблицы с номером главы/раздела
}{}

\AfterEndPreamble{
%% регистрируем счётчики в системе totcounter
    \regtotcounter{totalcount@figure}
    \regtotcounter{totalcount@table}       % Если иным способом поставить в преамбуле то ошибка в числе таблиц
    \regtotcounter{TotPages}               % Если иным способом поставить в преамбуле то ошибка в числе страниц
    \newtotcounter{totalappendix}
    \newtotcounter{totalchapter}

    % Применять запрет переносов только при печати оглавления:
    % Оборачиваем оригинальную \tableofcontents в локальную группу с высокими штрафами
    \let\origtableofcontents\tableofcontents
    \makeatletter
    \renewcommand{\tableofcontents}{%
        % \begingroup
            \hyphenpenalty=10000\relax
            % \exhyphenpenalty=10000\relax
    %         % Временно игнорируем все вызовы \addcontentsline внутри \tableofcontents,
    %         % чтобы не добавлять запись "Оглавление" в саму оглавление.
    %         % \let\addcontentsline\@gobblethree
            \exhyphenpenalty=10000\relax
            % зафиксировать межстрочный интервал и убрать лишние вертикальные отступы внутри оглавления
            % \OnehalfSpacing*
            \justifying
            \origtableofcontents
        % \endgroup
        % \endgroup
    }%
    \makeatother

    % Установить одинаковый левый отступ для section и более глубоких уровней в оглавлении.
    % Значение 2.55em совпадает с \setrmarg{1.55em} выбором отступа выше.
    \cftsetindents{section}{0.0em}{1.85em}
    \cftsetindents{subsection}{0.0em}{2.65em}
    \cftsetindents{subsubsection}{0.0em}{2.55em}
}

\makeatletter
\renewcommand*{\cftchapterformatpnum}[1]{%
  \hbox to \@pnumwidth{\cftdotfill{\cftchapterdotsep}\cftchapterpagefont#1}}
\renewcommand*{\cftsectionformatpnum}[1]{%
  \hbox to \@pnumwidth{\cftdotfill{\cftdotsep}\cftsectionpagefont#1}}
\renewcommand*{\cftsubsectionformatpnum}[1]{%
  \hbox to \@pnumwidth{\cftdotfill{\cftdotsep}\cftsubsectionpagefont#1}}
\renewcommand*{\cftsubsubsectionformatpnum}[1]{%
  \hbox to \@pnumwidth{\cftdotfill{\cftdotsep}\cftsubsubsectionpagefont#1}}
\makeatother

\newcommand{\unnumberedsection}[1]{%
  \section*{#1}%
  \addcontentsline{toc}{section}{#1}%
}
